\section*{附录A：关键配置与复现实验步骤}
本文实验通过YAML统一配置驱动，关键字段包括模型结构、补丁大小、优化器参数、风险任务特征筛选策略等。该方式便于复现实验并支持快速切换训练场景。

\subsection*{主要程序模块说明}
\begin{enumerate}
	\item 分割训练脚本：负责模型构建、损失计算、验证与检查点保存；
	\item 分割评估脚本：负责顺序补丁推理、体数据重建与NIfTI结果写出；
	\item 风险训练脚本：负责表格数据加载、交叉注意力模型训练、阈值搜索和测试评估；
	\item 数据模块：实现病例扫描、标签匹配、重采样、补丁采样与特征预处理。
\end{enumerate}

\subsection*{形态学特征说明}
系统可从分割掩膜中提取体积、表面积、长宽深、纵横比、球形度、紧致度等特征，作为风险预测候选输入。相关特征计算模块已在项目中实现，便于后续与临床统计分析结合。

\subsection*{项目复现实验步骤}
\begin{enumerate}
	\item 根据环境配置安装依赖并准备CTA影像与标签数据；
	\item 配置训练参数，运行分割训练脚本生成最优模型；
	\item 基于最优分割模型导出预测结果并准备风险表格数据；
	\item 运行风险训练脚本，得到测试集评估指标与报告文件。
\end{enumerate}
